\textbf{Elección y deducción del procedimiento.} El equilibrio tangencial exige $T=D$. Con la polar $D=qS(C_{D0}+KC_L^2)$ y $C_L=W/(qS)$, calcular sustentación, coeficiente de resistencia y fuerza en ese orden.

\begin{enumerate}

\item Formar presión dinámica.

\[\fitmath{q=\frac{\rho\cdot {V}^{2}}{2}}\]

\[\fitmath{q=\frac{(0.85)\cdot {(85)}^{2}}{2}\approx 3070.625\,\mathrm{Pa}}\]

\item Imponer equilibrio vertical.

\[\fitmath{C_L=\frac{W}{q\cdot S}}\]

\[\fitmath{C_L=\frac{(48000)}{(3070.625)\cdot (24)}\approx 0.6513332\,}\]

\item Evaluar la polar en el coeficiente requerido.

\[\fitmath{C_D=C_{D0}+K\cdot {C_L}^{2}}\]

\[\fitmath{C_D=(0.027)+(0.049)\cdot {(0.6513331976)}^{2}\approx 0.04778751\,}\]

\item Imponer equilibrio tangencial.

\[\fitmath{T=q\cdot S\cdot C_D}\]

\[\fitmath{T=(3070.625)\cdot (24)\cdot (0.04778751178)\approx 3521.701\,\mathrm{N}}\]

\end{enumerate}
